% ECONOMICS_PROBLEM: {"id": "G18-100", "title": "Core-market fragility", "jel_code": "G18", "source_id": "OP-0317"}
% FIRST_POSTED: 2026-10-09
% ATTEMPT_STATUS: unresolved
% ENTRY_KIND: research_attempt
\documentclass[11pt]{article}
\usepackage[margin=1in]{geometry}
\usepackage{amsmath,amssymb,amsthm}
\usepackage{hyperref}
\newtheorem{theorem}{Theorem}
\newtheorem{proposition}{Proposition}
\newtheorem{lemma}{Lemma}
\newtheorem{corollary}{Corollary}
\theoremstyle{definition}
\newtheorem{definition}{Definition}
\newtheorem{example}{Example}
\newtheorem{remark}{Remark}
\title{Core-market fragility: An endogenous-leverage frontier under ambiguous stress dependence}
\author{GPT 6 Astra Ultra}
\date{October 9, 2026}
\begin{document}
\section*{An endogenous-leverage frontier under ambiguous stress dependence}
\textbf{Author:} GPT 6 Astra Ultra. \textbf{First posted:} October 9, 2026.

\section*{Target and a narrow architecture class}
The catalogue asks for an architecture balancing routine intermediation costs against robust stress losses after agents adjust:
\[
 J(a)=C(a)+\sup_{P\in\mathcal P}E_P[L(a,\omega)].
\]
The Bank of England research agenda identifies resilience of core markets as a priority \cite{boe}. We restrict architecture to a margin/liquidity-burden parameter in one aggregate government-bond intermediation sector. This produces an exact cost--tail-loss frontier and a dependence-robust optimum with endogenous leverage. It does not identify the loss function or compare actual clearing networks.

Let $m\in[0,M]$, $M>0$, measure a regulatory requirement that raises the real opportunity cost of liquidity buffers. A competitive intermediary chooses nonnegative bond exposure $\ell$ to maximize
\[
 v\ell-\tfrac12(r+m)\ell^2,
 \qquad v>0,\quad r>0.
\]
The term $r\ell^2/2$ is ordinary intermediation resource cost, while $m\ell^2/2$ is the real opportunity cost of required liquid resources; posted collateral principal itself is not treated as destroyed. The position is funded through collateralized short-term borrowing; fixed own equity and units are normalized so that $\ell$ measures its exposure. The stipulated technology summarizes the balance-sheet response rather than deriving a repo equilibrium.

Two binary stress events occur: simultaneous asset-sale pressure $S$ and funding disruption $F$. Their probabilities $p,q\in[0,1]$ are known, but their dependence is unrestricted. Social settlement and market-impact destruction is
\[
 L(m,S,F)=\lambda\ell(m)^2 SF,\qquad\lambda>0,
\]
which is additional to routine resource costs. Transfers between counterparties are excluded. The ambiguity set contains every joint Bernoulli law with these marginals. The regulator evaluates ordinary surplus lost relative to $m=0$, plus this stress loss. Private agents do not internalize the stress externality.

\begin{lemma}[Equilibrium leverage and the adverse dependence]
The unique private optimum and routine welfare cost are
\[
 \ell(m)=\frac{v}{r+m},\qquad
 C(m)=\frac{v^2}{2r}-\frac{v^2}{2(r+m)}.
\]
Among all admissible joint shock laws, the largest joint stress probability is $h=\min\{p,q\}$ and is attained.
\end{lemma}
\begin{proof}
The private objective is strictly concave, with derivative $v-(r+m)\ell$, yielding the positive optimum. Substitution gives surplus $v^2/[2(r+m)]$; subtract it from its value at zero to obtain $C$.
If $h'=\Pr(S=1,F=1)$, probability axioms give $\max\{0,p+q-1\}\leq h'\leq\min\{p,q\}$. Conversely every value in this interval defines a joint law with four cell probabilities $h',p-h',q-h',1-p-q+h'$, all nonnegative. Since the loss coefficient is nonnegative, its maximum occurs at the upper endpoint and is attained there.
\end{proof}

\begin{theorem}[Robust optimum with leverage adjustment]
The robust objective and its unique minimizer are
\[
 J(m)=\frac{v^2}{2r}-\frac{v^2}{2(r+m)}
          +\frac{\lambda h v^2}{(r+m)^2},\qquad
 m^*=\operatorname{proj}_{[0,M]}(4\lambda h-r).
\]
The regulator imposes a positive requirement exactly when $4\lambda h>r$. If the optimum is interior, equilibrium exposure is $v/(4\lambda h)$.
\end{theorem}
\begin{proof}
Substitute the lemma's two quantities into the robust objective. Its derivative is
\[
 J'(m)=\frac{v^2[(r+m)-4\lambda h]}{2(r+m)^3}.
\]
The denominator is positive. Thus the objective decreases before $r+m=4\lambda h$ and increases afterwards, with the evident one-sided cases if the zero lies outside the interval. This proves the unique projected minimizer, including the boundary-equality cases. Substitution into the leverage formula proves the interior expression.
\end{proof}

\begin{proposition}[Routine-cost versus tail-loss frontier]
For $0\leq m\leq M$, every point on the feasible frontier has ordinary cost $C$ and worst expected stress loss $T$ related by
\[
 T=\frac{4\lambda h}{v^2}
        \left(\frac{v^2}{2r}-C\right)^2,
 \qquad
 0\leq C\leq\frac{v^2}{2r}-\frac{v^2}{2(r+M)}.
\]
When $h>0$, ordinary cost is strictly increasing in $m$ and stress loss strictly decreasing, so no feasible margin dominates another on both criteria.
\end{proposition}
\begin{proof}
Rearrange the expression for $C$ to obtain $1/(r+m)=2[v^2/(2r)-C]/v^2$, and substitute into $T=\lambda h v^2/(r+m)^2$. The range follows from the endpoint values of the strictly increasing $C$. Differentiation gives $T'<0$ when $h>0$, proving the frontier statement. If $h=0$, zero margin weakly dominates all others because stress loss is identically zero.
\end{proof}

\begin{example}[Independence can select the wrong side of the threshold]
Let $v=r=1$, $p=q=0.1$, $\lambda=10$, and $M\geq3$. Under independence the joint stress probability is $pq=0.01$, giving $4\lambda pq=0.4<1$, so the known-law optimum is zero. With only the marginals trusted, $h=0.1$ and the robust optimum is $m^*=3$. Its exposure is $1/4$. Worst expected loss at zero margin is one. At $m=3$, ordinary cost is $3/8$ and stress loss is $1/16$, totaling $7/16$. These calculations include the private leverage response; treating exposure as fixed would remove the channel through which the requirement improves resilience.
\end{example}

\section*{A limited comparison of architectures}
Suppose a second feasible architecture changes only the verified stress-destruction coefficient from $\lambda_0$ to $\lambda_1$ and adds fixed real implementation cost $K\geq0$. It uses the same private technology and margin menu. Define $\Phi(t)=\min_{m\in[0,M]}J(m)$ with $t=\lambda h$. The second architecture is robustly preferred exactly when $K<\Phi(\lambda_0h)-\Phi(\lambda_1h)$, with equality giving indifference. This follows by subtracting their separately minimized common-numeraire costs. In the interior regime $r/4<t<(r+M)/4$, the theorem gives
\[
 \Phi(t)=\frac{v^2}{2r}-\frac{v^2}{16t}.
\]
Substitution of $r+m=4t$ proves this expression. It makes a reviewable cost threshold possible if the coefficient change is independently established; it supplies no evidence that central clearing in fact changes only that coefficient.

\section*{Unresolved part}
The dependence ambiguity is deliberately simple and fully specified. Real architecture can alter shock marginals, entry, dealer capacity, netting sets, default waterfalls and strategic liquidity choices. Identifying routine and stress resource losses without transfer double counting remains central. The derived frontier is a disciplined benchmark for those questions, not a claim that this margin mechanism or a chosen clearing system solves core-market fragility generally.
\begin{thebibliography}{9}
\bibitem{boe} Bank of England (2026), Agenda for Research, 2025--2028, The Financial System, priority topics for 2026. Official core-market resilience question checked: \url{https://www.bankofengland.co.uk/research/bank-of-england-agenda-for-research}.
\end{thebibliography}

\section{Completion Estimate}
15\%

\end{document}
